% =====================================================================
%  Thème 4 — EXERCICES — Niveau MOYEN
% =====================================================================
\documentclass[11pt,a4paper]{article}
\usepackage[utf8]{inputenc}
\usepackage[T1]{fontenc}
\usepackage{lmodern}
\usepackage[french]{babel}
\usepackage{amsmath,amssymb,mathtools}
\usepackage{icomma}
\usepackage{xcolor}
\usepackage[most]{tcolorbox}
\usepackage{enumitem}
\usepackage{array}
\usepackage{booktabs}
\usepackage{fancyhdr}
\usepackage[hidelinks]{hyperref}
\usepackage[a4paper,margin=2.2cm,headheight=26pt,headsep=10pt]{geometry}

\definecolor{primary}{RGB}{223,87,99}
\definecolor{primarydark}{RGB}{150,42,55}

\tcbset{boxbase/.style={enhanced,breakable,boxrule=0.9pt,arc=2.5pt,
   left=3mm,right=3mm,top=2.3mm,bottom=2.3mm,
   fonttitle=\bfseries,coltitle=white,toptitle=0.6mm,bottomtitle=0.6mm}}
\newtcolorbox[auto counter]{exercice}[1][]{boxbase,colback=primary!4,colframe=primary,
  title={\textsc{Exercice}~\thetcbcounter\ifx\relax#1\relax\relax\else\textmd{ --- \itshape #1}\fi}}

\pagestyle{fancy}
\fancyhf{}
\fancyhead[L]{\small\textcolor{primarydark}{\textbf{Fiche 6} — Les limites}}
\fancyhead[R]{\small\textcolor{primarydark}{\textsc{Thème 4 — Moyen}}}
\fancyfoot[C]{\small\thepage}
\renewcommand{\headrulewidth}{0.8pt}
\renewcommand{\headrule}{\hbox to\headwidth{\color{primary}\leaders\hrule height \headrulewidth\hfill}}

\newcommand{\R}{\mathbb{R}}

\begin{document}

\begin{center}
{\Large\bfseries\color{primarydark}Thème 4 — Radicaux : facteur dominant \& conjugué}\\[2pt]
{\large\color{primary}Exercices — Niveau Moyen}
\end{center}
\vspace{4pt}
{\small\itshape On soigne le $|x|$ selon le côté, et on lève des indéterminations $\infty-\infty$ et
$\frac00$ par le conjugué.}
\vspace{6pt}

\begin{exercice}[le signe de $|x|$ change tout]
Soit $f(x)=\dfrac{\sqrt{x^2+1}}{x}$.
\begin{enumerate}[label=\alph*),leftmargin=2em,itemsep=3pt]
  \item Écrire $\sqrt{x^2+1}$ sous la forme $|x|\sqrt{1+\frac{1}{x^2}}$.
  \item Calculer $\displaystyle\lim_{x\to+\infty} f(x)$ (ici $|x|=x$).
  \item Calculer $\displaystyle\lim_{x\to-\infty} f(x)$ (ici $|x|=-x$).
\end{enumerate}
\end{exercice}

\begin{exercice}[différence de racines : $\infty-\infty$]
Calculer $\displaystyle\lim_{x\to+\infty}\bigl(\sqrt{x^2+x}-x\bigr)$.
\begin{enumerate}[label=\alph*),leftmargin=2em,itemsep=3pt]
  \item Vérifier que la forme directe est $\infty-\infty$.
  \item Multiplier par le conjugué $\sqrt{x^2+x}+x$ et simplifier le numérateur.
  \item Mettre $x$ en facteur au dénominateur et conclure.
\end{enumerate}
\end{exercice}

\begin{exercice}[$\frac00$ en un point]
Calculer $\displaystyle\lim_{x\to 3}\frac{\sqrt{x+1}-2}{x-3}$ par le conjugué du numérateur.
\end{exercice}

\begin{exercice}[somme de racine et polynôme]
Calculer $\displaystyle\lim_{x\to-\infty}\bigl(\sqrt{x^2+1}+x\bigr)$.
\begin{enumerate}[label=\alph*),leftmargin=2em,itemsep=3pt]
  \item Pourquoi la forme est-elle $\infty-\infty$ quand $x\to-\infty$ ? (penser au signe de $x$)
  \item Multiplier par le conjugué $\sqrt{x^2+1}-x$ et conclure.
\end{enumerate}
\end{exercice}

\begin{exercice}[racine sur quotient factorisable]
Calculer $\displaystyle\lim_{x\to 2}\frac{\sqrt{x+2}-2}{x^2-4}$.
\begin{enumerate}[label=\alph*),leftmargin=2em,itemsep=3pt]
  \item Factoriser le dénominateur.
  \item Multiplier par le conjugué du numérateur, simplifier le facteur commun et conclure.
\end{enumerate}
\end{exercice}

\end{document}
